% ============================================================================
%  PLANTILLA DE INFORME — Universidad Privada Domingo Savio
%  ----------------------------------------------------------------------------
%  Regla de oro del proyecto:
%    - formato/datos-institucion.sty -> datos compartidos (institución, logo)
%    - formato/preambulo.sty         -> toda la configuración (APA + Arial)
%    - formato/portada.tex           -> portada (datos institucionales)
%    - formato/referencias.bib       -> bibliografía BibTeX
%    - contenido/*.tex               -> el contenido real del informe
%    - anexos/*.tex                  -> material complementario
%    - main.tex                      -> integración + datos del documento
%
%  Compilar (desde la raíz de esta carpeta):
%    make
% ============================================================================

\documentclass[12pt,a4paper]{article}

\usepackage{formato/preambulo}
\usepackage{formato/datos-institucion}

% ============================================================================
%  DATOS DEL DOCUMENTO  (editar aquí en cada informe)
% ============================================================================
\newcommand{\tituloDocumento}{[Título del informe]}
\newcommand{\equipo}{[Nombre del equipo]}
\newcommand{\integrantes}{%
  [Integrante 1] \par
  [Integrante 2] \par
  [Integrante 3]}
\newcommand{\asignatura}{[Asignatura]}
\newcommand{\docente}{[Docente]}
\newcommand{\fechaEntrega}{DD/MM/AAAA}

\begin{document}

\pagenumbering{roman}
% ============================================================================
%  PORTADA  (formato APA 7.ª edición)
%  ----------------------------------------------------------------------------
%  Los datos provienen de las macros definidas en main.tex. Editar ahí.
% ============================================================================

\begin{titlepage}
  \thispagestyle{empty}
  \setstretch{1}
  \setlength{\parskip}{0pt}

  \begin{center}
    \vspace*{0.6cm}

    {\bfseries\fontsize{20}{24}\selectfont \universidad \par}
    {\bfseries\fontsize{18}{22}\selectfont \facultad \par}
    {\bfseries\fontsize{16}{19}\selectfont \carrera \par}

    \vspace{0.8cm}

    % Logo de la institución (colocar "logo-upds.png" u otro nombre en figuras/)
    \IfFileExists{figuras/logo.png}{%
      \includegraphics[width=0.30\textwidth]{figuras/logo.png}\par
    }{}

    \vspace{0.9cm}

    {\bfseries\Large \tituloDocumento \par}

    \vspace{1.0cm}

    {\large \integrantes \par}

    \vspace{1.0cm}

    {\normalsize \textbf{Asignatura:} \asignatura \par}
    {\normalsize \textbf{Docente:} \docente \par}
    {\normalsize \textbf{Fecha de entrega:} \fechaEntrega \par}

    \vspace{1.1cm}

    {\normalsize \textbf{Equipo:} \equipo \par}

    \vspace{0.9cm}

    {\normalsize \ciudad \par}

    \vspace{0.2cm}

    {\normalsize \the\year \par}
  \end{center}
\end{titlepage}

\newpage
{\indicesCompactos\tableofcontents}
\newpage
{\indicesCompactos\listoftables}
\newpage
{\indicesCompactos\listoffigures}
\clearpage

\pagenumbering{arabic}

% --- CAPÍTULO I --------------------------------------------------------------
\caratulaCapitulo{CAPÍTULO I}{INTRODUCCIÓN}
% ============================================================================
%  CAPÍTULO I: INTRODUCCIÓN
% ============================================================================

\section{Introducción}

[Escribir el contenido aquí.]


% --- CAPÍTULO II -------------------------------------------------------------
\caratulaCapitulo{CAPÍTULO II}{OBJETIVOS, ALCANCE Y LÍMITES}
% ============================================================================
%  CAPÍTULO II: OBJETIVOS, ALCANCE Y LÍMITES
% ============================================================================

\section{Objetivos del trabajo}

\subsection{Objetivo general}

[Describir el objetivo general del trabajo.]

\subsection{Objetivos específicos}

\begin{enumerate}
  \item Primer objetivo específico.
  \item Segundo objetivo específico.
  \item Tercer objetivo específico.
\end{enumerate}

\section{Tabla de ejemplo}

Tabla con el formato APA definido en el preámbulo (booktabs, sin líneas
verticales, interlineado sencillo):

\begin{table}[H]
  \centering
  \caption{Materiales y herramientas utilizados}
  \contenidoConFuente{%
    \begin{tablaAPA}
      \begin{tabular}{@{}p{3.6cm}p{8.2cm}p{3.2cm}@{}}
        \toprule
        \textbf{Recurso} & \textbf{Descripción} & \textbf{Versión} \\
        \midrule
        Debian GNU/Linux  & Sistema operativo base & 12 \\
        Nginx             & Servidor web           & 1.26 \\
        Samba AD DC       & Controlador de dominio & 4.19 \\
        \bottomrule
      \end{tabular}
    \end{tablaAPA}%
  }{Elaboración propia.}
\end{table}

% ============================================================================
%  CAPÍTULO II: OBJETIVOS, ALCANCE Y LÍMITES  (continúa)
% ============================================================================

\section{Alcance y Límites del Trabajo}

[Describir el alcance y los límites del trabajo.]

\subsection{Alcance}

[Escribir el alcance aquí.]

\subsection{Límites}

[Escribir los límites aquí.]

\section{Ecuación de ejemplo}

Una ecuación centrada con numeración:

\begin{equation}
  E = mc^2
\end{equation}


% --- CAPÍTULO III ------------------------------------------------------------
\caratulaCapitulo{CAPÍTULO III}{DESARROLLO DEL TRABAJO}
% ============================================================================
%  CAPÍTULO III: DESARROLLO DEL TRABAJO
% ============================================================================

\section{Desarrollo del trabajo}

\subsection{Figura de ejemplo}

Figura con título y línea FUENTE al estilo APA (usa el logo institucional
incluido en \texttt{figuras/logo.png}):

\figuraAPA{fig:logo}{Logo institucional UPDS}{\includegraphics[width=0.35\textwidth]{logo.png}}{Universidad Privada Domingo Savio.}

\subsection{Captura pendiente}

Esta macro muestra un recuadro mientras la captura no exista en
\texttt{figuras/}; cuando guardes la imagen con ese nombre, aparece
automáticamente:

\capturaAPA{fig:captura-ejemplo}{Ejemplo de captura de pantalla}{captura-ejemplo.png}{Pantalla de ejemplo del sistema.}

\subsection{Bloque de código}

\begin{lstlisting}[language=bash]
# Comandos de ejemplo
sudo apt update
sudo systemctl restart nginx
\end{lstlisting}

\subsection{Terminal}

\begin{lstlisting}[style=terminal]
$ systemctl status nginx
nginx.service - A high performance web server
   Active: active (running)
\end{lstlisting}

\subsection{Diagrama de arquitectura de red}

Ejemplo de topología cliente–router–servidor (como en los informes reales):

\figuraAPA{fig:red}{Arquitectura de red del laboratorio}{%
  \begin{tikzpicture}[
      nodo/.style={draw, rounded corners, minimum width=2.4cm,
                   minimum height=0.9cm, align=center, font=\small},
      nube/.style={draw, ellipse, align=center, font=\small, minimum width=2.2cm},
      linea/.style={-Stealth, thick}]
    \node[nodo] (pc) {Cliente};
    \node[nodo, right=1.8cm of pc] (rtr) {Router\textbackslash 192.168.0.1};
    \node[nube, right=1.8cm of rtr] (net) {Internet};
    \node[nodo, below=1.4cm of rtr] (srv) {Servidor Debian};
    \draw[linea] (pc) -- (rtr);
    \draw[linea] (rtr) -- (net);
    \draw[linea] (srv) -- (rtr);
  \end{tikzpicture}%
}{Elaboración propia.}

\subsection{Diagrama de estados (autómata)}

Ejemplo de máquina de estados finitos con transiciones etiquetadas:

\figuraAPA{fig:fsm}{Diagrama de estados de ejemplo}{%
  \begin{tikzpicture}[
      estado/.style={draw, circle, minimum size=1cm, font=\small},
      linea/.style={-Stealth, thick},
      nodo/.style={align=center}]
    \node[estado] (q0) {$q_0$};
    \node[estado, right=1.8cm of q0] (q1) {$q_1$};
    \node[estado, accepting, right=1.8cm of q1] (q2) {$q_2$};
    \draw[linea] (q0) -- node[above] {0} (q1);
    \draw[linea] (q1) -- node[above] {1} (q2);
    \draw[linea] (q1) to[bend left] node[above] {0} (q0);
    % Si se rompe la línea siguiente, el estado inicial desaparece (nodos y
    % posicionamiento relativos: siempre definir los nodos antes de usarlos).
  \end{tikzpicture}%
}{Elaboración propia.}

\subsection{Diagrama de flujo}

\figuraAPA{fig:flujo}{Flujo de decisión de ejemplo}{%
  \begin{tikzpicture}[
      paso/.style={draw, rounded corners, minimum width=2.6cm, minimum height=0.9cm, align=center, font=\small},
      decision/.style={draw, diamond, aspect=2, align=center, font=\small, inner sep=1pt},
      linea/.style={-Stealth, thick}]
    \node[paso] (ini) {Inicio};
    \node[decision, below=1cm of ini] (d) {¿Servicio activo?};
    \node[paso, below=1cm of d] (ok) {Continuar};
    \node[paso, right=1.6cm of d] (err) {Reiniciar};
    \draw[linea] (ini) -- (d);
    \draw[linea] (d) -- node[right, font=\scriptsize] {sí} (ok);
    \draw[linea] (d) -- node[above, font=\scriptsize] {no} (err);
  \end{tikzpicture}%
}{Elaboración propia.}

\subsection{Casos que "siempre daban error" (prueba del fix)}

El bloque siguiente usa flechas con la sintaxis moderna
(\texttt{->}, \texttt{>=stealth}) y texto entre nodos. **Antes** del fix
de babel (\texttt{\\shorthandoff\{<>\}} + librería \texttt{babel}) esto
fallaba con \texttt{Argument of \\language@active@arg> has an extra \}}.
Si compila, el fix funciona:

\figuraAPA{fig:fix}{Prueba del fix babel/TikZ}{%
  \begin{tikzpicture}[nodo/.style={draw, rounded corners, minimum width=2cm,
                     minimum height=0.8cm, align=center}]
    \node[nodo] (a) {Antes};
    \node[nodo, right=2cm of a] (b) {Después};
    \draw[->, >=stealth, thick] (a) -- node[above, font=\scriptsize]{tcp/80} (b);
  \end{tikzpicture}%
}{Elaboración propia.}

{\footnotesize\textbf{Errores que ya no deben aparecer:} tokens activos de
babel-spanish (\texttt{<}, \texttt{>}) rompiendo estilos de flecha;
texto de nodos con acentos dentro de argumentos de macro (usar
\texttt{\\shorthandoff} ya activo); depender de librerías no cargadas
(automata: cargar \texttt{automata} si se usa \texttt{accepting};
arrows.meta ya está cargado).\par}


% --- CONCLUSIONES ------------------------------------------------------------
\newpage
\encabezadoCapitulo{CONCLUSIONES}{}
\addcontentsline{toc}{part}{CONCLUSIONES}
\begingroup
\setcounter{capitulo}{4}
\setcounter{section}{0}
\renewcommand{\thesection}{4.}
\renewcommand{\theHsection}{4.}
\renewcommand{\thesubsection}{4.\arabic{subsection}}
\renewcommand{\theHsubsection}{4.\arabic{subsection}}
\renewcommand{\thesubsubsection}{4.\arabic{subsection}.\arabic{subsubsection}}
\renewcommand{\theHsubsubsection}{4.\arabic{subsection}.\arabic{subsubsection}}
% ============================================================================
%  CAPITULO IV: CONCLUSIONES
% ============================================================================

\section{Conclusiones}

\begin{enumerate}
  \item Primera conclusión del trabajo.
  \item Segunda conclusión del trabajo.
  \item Tercera conclusión del trabajo.
\end{enumerate}

\endgroup

% --- Referencias -------------------------------------------------------------
\newpage
\encabezadoCapitulo{REFERENCIAS}{}
\addcontentsline{toc}{part}{REFERENCIAS}
\bibliographystyle{apacite}
\nocite{*}
\bibliography{formato/referencias}

% --- Anexos -------------------------------------------------------------------
% \newpage
% \encabezadoCapitulo{ANEXOS}{}
% \pdfbookmark[0]{ANEXOS}{anexos}%
% \addcontentsline{toc}{part}{ANEXOS}
% \appendix
% % ============================================================================
%  ANEXO A: [TÍTULO DEL ANEXO]
%  Plantilla: copiar para cada anexo (anexo-a, anexo-b, ...).
% ============================================================================

\section{Anexo A: [T\'itulo del anexo]}

[Escribir el contenido del anexo aqu\'i.]

% Ejemplo de tabla:
% \begin{table}[H]
%   \centering
%   \caption{T\'itulo de la tabla}
%   \contenidoConFuente{%
%     \begin{tablaAPA}
%       \begin{tabular}{@{}p{3.6cm}p{8.2cm}p{3.2cm}@{}}
%         \toprule
%         \textbf{Columna 1} & \textbf{Columna 2} & \textbf{Columna 3} \\
%         \midrule
%         Dato 1 & Dato 2 & Dato 3 \\
%         \bottomrule
%       \end{tabular}
%     \end{tablaAPA}%
%   }{Elaboraci\'on propia.}
% \end{table}

\end{document}
